\documentclass[12pt]{article}
\usepackage[brazil]{babel}
\usepackage[utf8]{inputenc}
\usepackage[T1]{fontenc}
\usepackage{amsmath, amssymb}
\usepackage{geometry}
\usepackage{graphicx}
\usepackage{xcolor}

\definecolor{univBlue}{RGB}{0, 57, 145}
\definecolor{univLightBlue}{RGB}{71, 115, 169}
\definecolor{univSoftBlue}{RGB}{147, 172, 203}
\definecolor{maincolor}{RGB}{0, 57, 145}

\usepackage{fancyhdr}
\usepackage[most]{tcolorbox}

\newtcolorbox{block}[1]{%
  colback=univSoftBlue!10,
  colframe=maincolor,
  coltitle=white,
  title=#1,
  fonttitle=\bfseries,
  boxrule=1.5pt,
  arc=3mm,
  left=6pt,
  right=6pt,
  top=6pt,
  bottom=6pt
}

\geometry{margin=2.5cm}
\pagestyle{fancy}
\fancyhf{}
\usepackage{tikz}
\usetikzlibrary{positioning,arrows.meta,shapes.multipart}

\renewcommand{\headrulewidth}{2pt}
\renewcommand{\footrulewidth}{2pt}

\renewcommand{\headrule}{\hbox to\headwidth{%
  \color{maincolor}\leaders\hrule height \headrulewidth\hfill}}

\renewcommand{\footrule}{\hbox to\headwidth{%
  \color{maincolor}\leaders\hrule height \footrulewidth\hfill}}

\setlength{\headheight}{52pt}

\fancyhead[L]{%
  \includegraphics[height=1.2cm]{universidade_logo_novo.png}
}

\fancyhead[R]{%
  \raggedleft
  \textcolor{maincolor}{%
    \small
    \textbf{Lista de Exercícios — Herança e Composição}\\
    Me.~Sergio Souza Novak\\
    \textit{Faculdade de Engenharia de Software}
  }
}

\fancyfoot[C]{%
  \textcolor{maincolor}{\thepage}
}

% Estilo reutilizável para diagramas UML
\tikzset{
  uml/.style={
    draw=maincolor,
    rectangle split,
    rectangle split parts=3,
    fill=univSoftBlue!20,
    minimum width=5cm,
    font=\small\ttfamily,
    align=left
  },
  heranca/.style={-{Triangle[open]}, thick, maincolor},
  composicao/.style={-{Diamond[fill=maincolor]}, thick, maincolor}
}

\begin{document}
% ============================================================
Lista de implementação de Herança e Composição. Implemente na linguagem Java todos os exercícios propostos.
\section*{Formato de Entrega}


\begin{block}{Instruções para Submissão}
A entrega desta lista de exercícios deve ser feita através de um \textbf{único arquivo .zip}.
Dentro deste arquivo, deve haver uma pasta estruturada para cada exercício (ex: \texttt{Exercicio1}, \texttt{Exercicio2}, ..., \texttt{Exercicio8}).

\textbf{Atenção:} Certifique-se de que cada pasta contenha \textbf{todos} os arquivos necessários (\texttt{.java}) para a correta compilação e execução do respectivo exercício.
\end{block}


% ============================================================
\section*{1) Herança — Veículos de Transporte}
Implemente uma hierarquia de classes para representar veículos de transporte. A superclasse \texttt{Veiculo} possui atributos e métodos comuns, e as subclasses especializam o comportamento.

\begin{center}
\begin{tikzpicture}[node distance=1.8cm and 0cm]

\node[uml] (veiculo) {
    \textbf{Veiculo}
    \nodepart{second}
    \# placa : String \\
    \# velocidadeMax : double \\
    \# kmRodados : double
    \nodepart{third}
    + getPlaca() : String \\
    + setPlaca(p:String) : void \\
    + getVelocidadeMax() : double \\
    + setVelocidadeMax(v:double) : void \\
    + getKmRodados() : double \\
    + setKmRodados(k:double) : void \\
    + acelerar(km:double) : void
};

\node[uml, below left=of veiculo, xshift=2cm] (carro) {
    \textbf{Carro}
    \nodepart{second}
    - numPassageiros : int
    \nodepart{third}
    + getNumPassageiros() : int \\
    + setNumPassageiros(n:int) : void \\
    + acelerar(km:double) : void
};

\node[uml, below right=of veiculo, xshift=-1.5cm] (caminhao) {
    \textbf{Caminhao}
    \nodepart{second}
    - cargaMaxTon : double
    \nodepart{third}
    + getCargaMaxTon() : double \\
    + setCargaMaxTon(c:double) : void \\
    + acelerar(km:double) : void
};

\draw[heranca] (carro.north) -- ++(0,0.5) -| (veiculo.south);
\draw[heranca] (caminhao.north) -- ++(0,0.5) -| (veiculo.south);

\end{tikzpicture}
\end{center}

\subsection*{Requisitos}
\begin{itemize}
    \item A classe \texttt{Veiculo} deve possuir o método \texttt{acelerar(km)} que incrementa \texttt{kmRodados}.
    \item \texttt{Carro} sobrescreve \texttt{acelerar()} usando \texttt{super} e adiciona uma verificação de limite de velocidade.
    \item \texttt{Caminhao} sobrescreve \texttt{acelerar()} reduzindo a eficiência em 30\% se a carga estiver acima de 80\% da capacidade.
    \item Crie uma classe \texttt{Main} que instancie ambos os veículos e teste os métodos.
\end{itemize}

% ============================================================
\section*{2) Herança — Funcionários de uma Empresa}
Modele um sistema de folha de pagamento com uma superclasse \texttt{Funcionario} e duas subclasses que calculam o salário de formas diferentes.

\begin{center}
\begin{tikzpicture}[node distance=1.8cm and -0.5cm]

\node[uml] (func) {
    \textbf{Funcionario}
    \nodepart{second}
    \# nome : String \\
    \# salarioBase : double
    \nodepart{third}
    + getNome() : String \\
    + setNome(n:String) : void \\
    + getSalarioBase() : double \\
    + setSalarioBase(s:double) : void \\
    + calcularSalario() : double
};

\node[uml, below left=of func, xshift=2cm] (gerente) {
    \textbf{Gerente}
    \nodepart{second}
    - bonus : double
    \nodepart{third}
    + getBonus() : double \\
    + setBonus(b:double) : void \\
    + calcularSalario() : double
};

\node[uml, below right=of func, xshift=0.5cm] (estagiario) {
    \textbf{Estagiario}
    \nodepart{second}
    - horasTrabalhadas : int
    \nodepart{third}
    + getHorasTrabalhadas() : int \\
    + setHorasTrabalhadas(h:int) : void \\
    + calcularSalario() : double
};

\draw[heranca] (gerente.north) -- ++(0,0.5) -| (func.south);
\draw[heranca] (estagiario.north) -- ++(0,0.5) -| (func.south);

\end{tikzpicture}
\end{center}

\subsection*{Requisitos}
\begin{itemize}
    \item \texttt{Gerente} calcula o salário como \texttt{salarioBase + bonus}.
    \item \texttt{Estagiario} calcula o salário com base nas horas trabalhadas multiplicadas por um valor fixo por hora (ex: R\$ 15,00).
    \item Utilize \texttt{super()} para inicializar os atributos da superclasse nos construtores.
    \item Crie um \texttt{Main} que imprima o salário de cada tipo de funcionário.
\end{itemize}

% ============================================================
\section*{3) Herança — Equipamentos Médicos}
Modele uma hierarquia de classes de equipamentos médicos de raio-X. A superclasse define o comportamento básico de emissão de radiação.

\begin{center}
\begin{tikzpicture}[node distance=1.8cm and -1cm]

\node[uml] (raiox) {
    \textbf{EquipamentoRaioX}
    \nodepart{second}
    \# modelo : String \\
    \# fabricante : String
    \nodepart{third}
    + getModelo() : String \\
    + setModelo(m:String) : void \\
    + getFabricante() : String \\
    + setFabricante(f:String) : void \\
    + ligar() : void \\
    + emitirRadiacao() : void
};

\node[uml, below left=of raiox, xshift=-1cm] (tomo) {
    \textbf{Tomografo}
    \nodepart{second}
    - numCortes : int
    \nodepart{third}
    + getNumCortes() : int \\
    + setNumCortes(n:int) : void \\
    + emitirRadiacao() : void \\
    + reconstruir3D() : void
};

\node[uml, below=of raiox] (mamo) {
    \textbf{Mamografo}
    \nodepart{second}
    - pressao : double
    \nodepart{third}
    + getPressao() : double \\
    + setPressao(p:double) : void \\
    + emitirRadiacao() : void \\
    + comprimir() : void
};

\node[uml, below right=of raiox, xshift=1cm] (densito) {
    \textbf{Densitometro}
    \nodepart{second}
    \nodepart{third}
    + emitirRadiacao() : void \\
    + calcularTScore() : double
};

\draw[heranca] (tomo.north) -- ++(0,0.5) -| (raiox.south);
\draw[heranca] (mamo.north) -- (raiox.south);
\draw[heranca] (densito.north) -- ++(0,0.5) -| (raiox.south);

\end{tikzpicture}
\end{center}

\subsection*{Requisitos}
\begin{itemize}
    \item Cada subclasse deve sobrescrever \texttt{emitirRadiacao()} chamando \texttt{super.emitirRadiacao()} e adicionando comportamento específico.
    \item \texttt{Tomografo} imprime a quantidade de cortes realizados.
    \item \texttt{Mamografo} imprime a pressão de compressão aplicada.
    \item \texttt{Densitometro} calcula e retorna um T-Score fictício.
    \item Teste todos os equipamentos em um \texttt{Main}.
\end{itemize}

% ============================================================
\section*{4) Composição — Computador e seus Componentes}
Implemente um sistema onde um \texttt{Computador} é composto por um \texttt{Processador} e uma \texttt{MemoriaRAM}. Sem o computador, as partes não existem.

\begin{center}
\begin{tikzpicture}[node distance=2cm and -2.5cm]

\node[uml] (pc) {
    \textbf{Computador}
    \nodepart{second}
    - marca : String \\
    - processador : Processador \\
    - memoria : MemoriaRAM
    \nodepart{third}
    + getMarca() : String \\
    + setMarca(m:String) : void \\
    + getProcessador() : Processador \\
    + setProcessador(p:Processador) : void \\
    + getMemoria() : MemoriaRAM \\
    + setMemoria(m:MemoriaRAM) : void \\
    + ligar() : void \\
    + desligar() : void
};

\node[uml, below left=of pc] (proc) {
    \textbf{Processador}
    \nodepart{second}
    - modelo : String \\
    - nucleos : int \\
    - frequenciaGHz : double
    \nodepart{third}
    + getModelo() : String \\
    + setModelo(m:String) : void \\
    + getNucleos() : int \\
    + setNucleos(n:int) : void \\
    + getFrequenciaGHz() : double \\
    + setFrequenciaGHz(f:double) : void \\
    + processar() : void
};

\node[uml, below right=of pc] (ram) {
    \textbf{MemoriaRAM}
    \nodepart{second}
    - capacidadeGB : int \\
    - tipo : String
    \nodepart{third}
    + getCapacidadeGB() : int \\
    + setCapacidadeGB(c:int) : void \\
    + getTipo() : String \\
    + setTipo(t:String) : void \\
    + armazenar() : void
};

\draw[composicao] (proc.north) -- ++(0,0.5) -| (pc.south);
\draw[composicao] (ram.north) -- ++(0,0.5) -| (pc.south);

\end{tikzpicture}
\end{center}

\subsection*{Requisitos}
\begin{itemize}
    \item O \texttt{Computador} deve criar os objetos \texttt{Processador} e \texttt{MemoriaRAM} internamente (composição forte).
    \item O método \texttt{ligar()} deve delegar para \texttt{processador.processar()} e \texttt{memoria.armazenar()}.
    \item Todos os atributos devem ser privados com acesso via métodos.
    \item Crie um \texttt{Main} demonstrando a criação e o uso do computador.
\end{itemize}

% ============================================================
\section*{5) Composição — Casa Inteligente}
Modele uma \texttt{CasaInteligente} composta por um \texttt{SistemaSeguranca} e um \texttt{SistemaIluminacao}. A casa gerencia os subsistemas.

\begin{center}
\begin{tikzpicture}[node distance=2cm and -4cm]

\node[uml] (casa) {
    \textbf{CasaInteligente}
    \nodepart{second}
    - endereco : String \\
    - seguranca : SistemaSeguranca \\
    - iluminacao : SistemaIluminacao
    \nodepart{third}
    + getEndereco() : String \\
    + setEndereco(e:String) : void \\
    + getSeguranca() : SistemaSeguranca \\
    + setSeguranca(s:SistemaSeguranca) : void \\
    + getIluminacao() : SistemaIluminacao \\
    + setIluminacao(i:SistemaIluminacao) : void \\
    + ativarModoNoturno() : void \\
    + ativarModoAusente() : void
};

\node[uml, below left=of casa] (seg) {
    \textbf{SistemaSeguranca}
    \nodepart{second}
    - ativo : boolean \\
    - alarmeDisparado : boolean
    \nodepart{third}
    + isAtivo() : boolean \\
    + setAtivo(a:boolean) : void \\
    + isAlarmeDisparado() : boolean \\
    + setAlarmeDisparado(a:boolean) : void \\
    + ativar() : void \\
    + desativar() : void \\
    + dispararAlarme() : void
};

\node[uml, below right=of casa] (luz) {
    \textbf{SistemaIluminacao}
    \nodepart{second}
    - intensidade : int \\
    - ligado : boolean
    \nodepart{third}
    + getIntensidade() : int \\
    + setIntensidade(i:int) : void \\
    + isLigado() : boolean \\
    + setLigado(l:boolean) : void \\
    + ligar() : void \\
    + desligar() : void \\
    + ajustarIntensidade(v:int) : void
};

\draw[composicao] (seg.north) -- ++(0,0.5) -| (casa.south);
\draw[composicao] (luz.north) -- ++(0,0.5) -| (casa.south);

\end{tikzpicture}
\end{center}

\subsection*{Requisitos}
\begin{itemize}
    \item \texttt{ativarModoNoturno()} deve ligar a segurança e reduzir a iluminação para 20\%.
    \item \texttt{ativarModoAusente()} deve ativar a segurança e desligar a iluminação.
    \item Demonstre a delegação de responsabilidades no \texttt{Main}.
\end{itemize}

% ============================================================
\section*{6) Herança + Composição — Universidade}
Combine herança e composição. Uma \texttt{Universidade} possui (composição) vários \texttt{Curso}s. \texttt{Curso} é superclasse de \texttt{CursoPresencial} e \texttt{CursoEAD}.

\begin{center}
\begin{tikzpicture}[node distance=1.8cm and -2cm]

\node[uml] (univ) {
    \textbf{Universidade}
    \nodepart{second}
    - nome : String \\
    - cursos : Curso[]
    \nodepart{third}
    + getNome() : String \\
    + setNome(n:String) : void \\
    + getCursos() : Curso[] \\
    + setCursos(c:Curso[]) : void \\
    + adicionarCurso(c:Curso) : void \\
    + listarCursos() : void
};

\node[uml, below=2.5cm of univ] (curso) {
    \textbf{Curso}
    \nodepart{second}
    \# nome : String \\
    \# duracaoSemestres : int
    \nodepart{third}
    + getNome() : String \\
    + setNome(n:String) : void \\
    + getDuracaoSemestres() : int \\
    + setDuracaoSemestres(d:int) : void \\
    + getInfo() : String
};

\node[uml, below left=of curso, xshift=-0.5cm] (presencial) {
    \textbf{CursoPresencial}
    \nodepart{second}
    - campus : String
    \nodepart{third}
    + getCampus() : String \\
    + setCampus(c:String) : void \\
    + getInfo() : String
};

\node[uml, below right=of curso, xshift=0.5cm] (ead) {
    \textbf{CursoEAD}
    \nodepart{second}
    - plataforma : String
    \nodepart{third}
    + getPlataforma() : String \\
    + setPlataforma(p:String) : void \\
    + getInfo() : String
};


\draw[composicao] (curso.north) -- (univ.south);
\draw[heranca] (presencial.north) -- ++(0,0.5) -| (curso.south);
\draw[heranca] (ead.north) -- ++(0,0.5) -| (curso.south);

\end{tikzpicture}
\end{center}

\subsection*{Requisitos}
\begin{itemize}
    \item \texttt{CursoPresencial} inclui o campus no \texttt{getInfo()}.
    \item \texttt{CursoEAD} inclui a plataforma no \texttt{getInfo()}.
    \item A \texttt{Universidade} deve ser capaz de listar todos os cursos com suas informações específicas.
    \item Utilize \texttt{super.getInfo()} para reaproveitar a lógica da superclasse.
\end{itemize}

% ============================================================
\section*{7) Herança — Contas Bancárias}
Modele um sistema bancário com uma superclasse \texttt{ContaBancaria} e duas subclasses com regras de saque distintas.

\begin{center}
\begin{tikzpicture}[node distance=1.8cm and -1.5cm]

\node[uml] (conta) {
    \textbf{ContaBancaria}
    \nodepart{second}
    \# titular : String \\
    \# saldo : double
    \nodepart{third}
    + getTitular() : String \\
    + setTitular(t:String) : void \\
    + getSaldo() : double \\
    + setSaldo(s:double) : void \\
    + depositar(valor:double) : void \\
    + sacar(valor:double) : boolean \\
    + getSaldo() : double
};

\node[uml, below left=of conta, xshift=2cm] (corrente) {
    \textbf{ContaCorrente}
    \nodepart{second}
    - limiteEspecial : double
    \nodepart{third}
    + getLimiteEspecial() : double \\
    + setLimiteEspecial(l:double) : void \\
    + sacar(valor:double) : boolean
};

\node[uml, below right=of conta, xshift=-1.5cm] (poupanca) {
    \textbf{ContaPoupanca}
    \nodepart{second}
    - taxaRendimento : double
    \nodepart{third}
    + getTaxaRendimento() : double \\
    + setTaxaRendimento(t:double) : void \\
    + sacar(valor:double) : boolean \\
    + aplicarRendimento() : void
};

\draw[heranca] (corrente.north) -- ++(0,0.5) -| (conta.south);
\draw[heranca] (poupanca.north) -- ++(0,0.5) -| (conta.south);

\end{tikzpicture}
\end{center}

\subsection*{Requisitos}
\begin{itemize}
    \item \texttt{ContaCorrente} permite saque até \texttt{saldo + limiteEspecial}.
    \item \texttt{ContaPoupanca} só permite saque se o saldo permanecer positivo (sem limite especial).
    \item O método \texttt{aplicarRendimento()} incrementa o saldo com base na \texttt{taxaRendimento}.
    \item Utilize \texttt{super()} nos construtores e \texttt{super.sacar()} onde fizer sentido.
\end{itemize}

% ============================================================
\section*{8) Composição — Pedido de Restaurante}
Modele um sistema de pedidos onde um \texttt{Pedido} é composto por vários \texttt{ItemPedido}s. Cada item referencia um \texttt{Produto}.

\begin{center}
\begin{tikzpicture}[node distance=2cm and -1cm]

\node[uml] (pedido) {
    \textbf{Pedido}
    \nodepart{second}
    - numero : int \\
    - itens : ItemPedido[] \\
    - cliente : String
    \nodepart{third}
    + getNumero() : int \\
    + setNumero(n:int) : void \\
    + getItens() : ItemPedido[] \\
    + setItens(i:ItemPedido[]) : void \\
    + getCliente() : String \\
    + setCliente(c:String) : void \\
    + adicionarItem(i:ItemPedido) : void \\
    + calcularTotal() : double \\
    + imprimirResumo() : void
};

\node[uml, below left=of pedido, xshift=1.5cm] (item) {
    \textbf{ItemPedido}
    \nodepart{second}
    - produto : Produto \\
    - quantidade : int
    \nodepart{third}
    + getProduto() : Produto \\
    + setProduto(p:Produto) : void \\
    + getQuantidade() : int \\
    + setQuantidade(q:int) : void \\
    + getSubtotal() : double
};

\node[uml, below right=of pedido, xshift=-0.5cm] (produto) {
    \textbf{Produto}
    \nodepart{second}
    - nome : String \\
    - preco : double
    \nodepart{third}
    + getNome() : String \\
    + setNome(n:String) : void \\
    + getPreco() : double \\
    + setPreco(p:double) : void
};

\draw[composicao] (item.north) -- ++(0,0.5) -| (pedido.south);
\draw[composicao] (produto.north) -- ++(0,0.5) -| (pedido.south);

\end{tikzpicture}
\end{center}

\subsection*{Requisitos}
\begin{itemize}
    \item O \texttt{Pedido} gerencia a lista de itens e calcula o total somando os subtotais.
    \item \texttt{ItemPedido} calcula o subtotal como \texttt{produto.getPreco() * quantidade}.
    \item \texttt{imprimirResumo()} deve listar todos os itens com nome, quantidade, subtotal e o total geral.
    \item Crie um \texttt{Main} que monte um pedido com pelo menos 3 itens e imprima o resumo.
\end{itemize}


\end{document}
